%%%%%%%%%%%%%%%%
%%% gen 37 %%%
%%%%%%%%%%%%%%%%


\Chapter{37}


\Verse{1}
{
	\hP{וַיֵּשֶׁב יַעֲקֹב בְּאֶרֶץ מְגוּרֵי אָבִיו בְּאֶרֶץ כְּנָעַן׃}
}
{
	\eP{
		And Jacob lived in the land of his father’s residences, in
		the land of Canaan.
	}
}
{
	Jacob lived near dad's house.
}


\Verse{2}
{
	\hR{אֵלֶּה תֹּלְדוֹת יַעֲקֹב}
	\hJ{יוֹסֵף בֶּן שְׁבַע עֶשְׂרֵה שָׁנָה הָיָה רֹעֶה אֶת אֶחָיו בַּצֹּאן וְהוּא נַעַר אֶת בְּנֵי בִלְהָה וְאֶת בְּנֵי זִלְפָּה נְשֵׁי אָבִיו וַיָּבֵא יוֹסֵף אֶת דִּבָּתָם רָעָה אֶל אֲבִיהֶם׃}
}
{
	\eR{
		These are the records of Jacob:
	}
	\eJ{
		Joseph, at seventeen years old, had been tending the sheep
		with his brothers, and he was a boy with the sons of Bilhah
		and the sons of Zilpah, his father’s wives; and Joseph
		brought a bad report of them to their father.
	}
}
{
	These are the records of Jacob.

	\aB{Oh Christ not another begat list.}

	Joseph was seventeen and watching the sheep with his brothers.

	\aB{Oh thank god it's a story.}

	Joseph goes to dad and complains about the older bros.
}


\Verse{3}
{
	\hE{וְיִשְׂרָאֵל אָהַב אֶת יוֹסֵף מִכּל בָּנָיו כִּי בֶן זְקֻנִים הוּא לוֹ}
	\hJ{וְעָשָׂה לוֹ כְּתֹנֶת פַּסִּים׃}
}
{
	\eE{
		And Israel had loved Joseph the most of all his children
		because he was a son of old age to him.
	}
	\eJ{
		And he made him a coat of many colors.
	}
}
{
	Joseph was Israel's favorite because obviously, Rachel.

	\aB{So he gives him a really rad shirt.}%
	\fC{Nobody knows what the Hebrew word here means.}
	\fB{Here the word is PSYm, it's YSP (Joseph) backwards.}
	\fC{Joseph (YSwP) has an internal waw.}
	\fB{Cross and Freedman point out that early Hebrew basically didn’t 
		do internal \textit{matres lectionis} yet. So if J was written
		around the ninth century, most of those helpful little \textit{waw}s
		and \textit{yod}s telling you what vowel to read just weren’t there.
		The original spelling would’ve been way more stripped-down and
		ambiguous. F. M. Cross and D. N. Freedman,
		\textit{Early Hebrew Orthography} (1952).
	}
	\fC{Red is learning quick.}

%	\footnote{
%		coat of many colors. We have no idea what the Hebrew means, so I have
%		retained the traditional “many colors.” Some have tried to derive
%		something about its meaning from the story of Tamar, who is wearing one
%		when her brother Amnon rapes her (2 Samuel 13). But this is difficult
%		because these two stories appear to be intentionally connected. (I
%		brought evidence that they come from the same source in The Hidden Book
%		in the Bible.) It is hardly coincidental that the two people who wear a
%		coat of many colors in the Bible are both victims of violence by their
%		brothers, and that both coats are torn. The significance of the coat,
%		therefore, is as a symbol of injustice among siblings. This significance
%		is doubled when we see the fate of the coat. See the comment on 37:31.
%	}
}


\Verse{4}
{
	\hE{וַיִּרְאוּ אֶחָיו כִּי אֹתוֹ אָהַב אֲבִיהֶם מִכּל אֶחָיו וַיִּשְׂנְאוּ אֹתוֹ וְלֹא יָכְלוּ דַּבְּרוֹ לְשָׁלֹם׃}
}
{
	\eE{
		And his brothers saw that their father loved him the most
		of all his brothers. And they hated him. And they were not
		able to speak a greeting to him.
	}
}
{
	His brothers are jealous and hate him.
%	\footnote{
%		to speak a greeting to him. Literally “to speak to him of peace,” it
%		appears to refer to the standard Hebrew greetings: one says šlôm
%		(peace/well-being) for hello; and one says hšlôm lô (does he have
%		peace?) for “How is he?” (Gen 29:6). That is, the brothers cannot even
%		bid him “hello” or ask how he is. This is rendered ironic by the fact
%		that Joseph’s father will soon send him to check on his brothers,
%		telling him, “See how your brothers are [šlôm].” It is rendered doubly
%		ironic later when Joseph tells the Pharaoh that it is possible to
%		interpret Pharaoh’s dream because “God will answer regarding Pharaoh’s
%		well-being [šlôm],” suggesting that God shows even more care for the
%		Pharaoh than Joseph’s brothers are able to show him. (And see the
%		comment on 43:27.)
%	}
}
\Verse{5}
{
	\hJ{וַיַּחֲלֹם יוֹסֵף חֲלוֹם וַיַּגֵּד לְאֶחָיו וַיּוֹסִפוּ עוֹד שְׂנֹא אֹתוֹ׃}
}
{
	\eJ{
		And Joseph had a dream and told it to his brothers, and
		they went on to hate him more.
	}
}
{
	Joseph has a dream.
	
	It made them hate him more.\fA{Lit. \heb{יּוֹסִפוּ} in the original.}
}

\Verse{6}
{
	\hJ{וַיֹּאמֶר אֲלֵיהֶם שִׁמְעוּ נָא הַחֲלוֹם הַזֶּה אֲשֶׁר חָלָמְתִּי׃}
}
{
	\eJ{And he said to them, “Listen to this dream that I had:}
}
{
	So he's like ``Hey guys, listen to this dream I had.''
}

\Verse{7}
{
	\hJ{וְהִנֵּה אֲנַחְנוּ מְאַלְּמִים אֲלֻמִּים בְּתוֹךְ הַשָּׂדֶה וְהִנֵּה קָמָה אֲלֻמָּתִי וְגַם נִצָּבָה וְהִנֵּה תְסֻבֶּינָה אֲלֻמֹּתֵיכֶם וַתִּשְׁתַּחֲוֶיןָ לַאֲלֻמָּתִי׃}
}
{
	\eJ{
		and here we were binding sheaves in the field; and here
		was my sheaf rising and standing up, too; and, here, your
		sheaves surrounded and bowed to my sheaf.”
	}
}
{
	We were out in the field tying up grain bundles.
	
	And mine starts standing up all straight and long
	and you guys' bundles all surround it and start
	bending over ha.
}

\Verse{8}
{
	\hJ{וַיֹּאמְרוּ לוֹ אֶחָיו הֲמָלֹךְ תִּמְלֹךְ עָלֵינוּ אִם מָשׁוֹל תִּמְשֹׁל בָּנוּ וַיּוֹסִפוּ עוֹד שְׂנֹא אֹתוֹ עַל חֲלֹמֹתָיו וְעַל דְּבָרָיו׃}
}
{
	\eJ{
		And his brothers said to him, “Will you rule over us?!
		Will you dominate us?!” And they went on to hate him more
		because of his dreams and because of his words.
	}
}
{
	So his brothers are like ``What the hell kind of a dream is \textsl{that}?!''

	Cuz obviously.

	They hate him even more now.
}

\Verse{9}
{
	\hJ{וַיַּחֲלֹם עוֹד חֲלוֹם אַחֵר וַיְסַפֵּר אֹתוֹ לְאֶחָיו וַיֹּאמֶר הִנֵּה חָלַמְתִּי חֲלוֹם עוֹד וְהִנֵּה הַשֶּׁמֶשׁ וְהַיָּרֵחַ וְאַחַד עָשָׂר כּוֹכָבִים מִשְׁתַּחֲוִים לִי׃}
}
{
	\eJ{
		And he had yet another dream and told it to his brothers.
		And he said, “Here, I’ve had another dream, and here were
		the sun and the moon and eleven stars bowing to me.”
	}
}
{
	So he has another dream and tells them again.

	He's like ``The sun and moon and eleven stars start bowing to me.''

}

\Verse{10}
{
	\hJ{וַיְסַפֵּר אֶל אָבִיו וְאֶל אֶחָיו וַיִּגְעַר בּוֹ אָבִיו וַיֹּאמֶר לוֹ מָה הַחֲלוֹם הַזֶּה אֲשֶׁר חָלָמְתָּ הֲבוֹא נָבוֹא אֲנִי וְאִמְּךָ וְאַחֶיךָ לְהִשְׁתַּחֲוֺת לְךָ אָרְצָה׃}
}
{
	\eJ{
		And he told it to his father and to his brothers. And his
		father was annoyed at him and said to him, “What is this
		dream that you’ve had? Shall we come, I and your mother
		and your brothers, to bow to you to the ground?!”
	}
}
{
	So he tells his older brothers and dad.

	All eleven of them, plus dad.

	And his dad is like ``What the hell kind of a dream is that?''

	Somehow everyone sees through his highly obscure non-obvious imagery.	
%	\footnote{
%		What is this dream that you’ve had? No one seems to know that the dreams
%		are prophetic: not the brothers, not Jacob, not even Joseph himself.
%		Joseph will be able to interpret other people’s dreams—but not his own.
%		Like many of us, he will learn what every psychoanalyst knows: that
%		seeing the meaning of one’s own dreams is the hardest.
%	}
}

\Verse{11}
{
	\hJ{וַיְקַנְאוּ בוֹ אֶחָיו וְאָבִיו שָׁמַר אֶת הַדָּבָר׃}
}
{
	\eJ{
		And his brothers were jealous of him, and his father took
		note of the thing.
	}
}
{
	His brothers keep being jealous.
	
	His dad is like ``Noted.''
}

\Verse{12}
{
	\hE{וַיֵּלְכוּ אֶחָיו לִרְעוֹת אֶת צֹאן אֲבִיהֶם בִּשְׁכֶם׃}
}
{
	\eE{
		And his brothers went to feed their father’s sheep in
		Shechem.
	}
}
{
	His brothers go to feed dad's sheep in that town they genocided
	with penis tricks.
}

\Verse{13}
{
	\hE{וַיֹּאמֶר יִשְׂרָאֵל אֶל יוֹסֵף הֲלוֹא אַחֶיךָ רֹעִים בִּשְׁכֶם לְכָה וְאֶשְׁלָחֲךָ אֲלֵיהֶם וַיֹּאמֶר לוֹ הִנֵּנִי׃}
}
{
	\eE{
		And Israel said to Joseph, “Aren’t your brothers feeding
		in Shechem? Come on and I’ll send you to them.” And he
		said to him, “I’m here.”
	}
}
{
	So Israel's like, ``Hey, your brothers are over in that town \textsl{they}
	genocided. Get out of here and go over there.''
	
	And he's like ``Whatcha need pops?''
}

\Verse{14}
{
	\hE{וַיֹּאמֶר לוֹ לֶךְ נָא רְאֵה אֶת שְׁלוֹם אַחֶיךָ וְאֶת שְׁלוֹם הַצֹּאן וַהֲשִׁבֵנִי דָּבָר וַיִּשְׁלָחֵהוּ מֵעֵמֶק חֶבְרוֹן וַיָּבֹא שְׁכֶמָה׃}
}
{
	\eE{
		And he said to him, “Go, see how your brothers are and how
		the sheep are and bring me back word.” And he sent him
		from the valley of Hebron. And he came to Shechem.
	}
}
{
	So dad's like ``Go check on your brothers and the sheep and let me know.''

	So Joseph heads to Shechem.
	
	The town, not the guy.
	
	The guy is dead.
}

\Verse{15}
{
	\hE{וַיִּמְצָאֵהוּ אִישׁ וְהִנֵּה תֹעֶה בַּשָּׂדֶה וַיִּשְׁאָלֵהוּ הָאִישׁ לֵאמֹר מַה תְּבַקֵּשׁ׃}
}
{
	\eE{
		And a man found him, and here he was straying in a field.
		And the man asked him, saying, “What are you looking for?”
	}
}
{
	A man finds him. He was doing something out in a field.
	
	The guy's like ``Whatcha doing?''
}

\Verse{16}
{
	\hE{וַיֹּאמֶר אֶת אַחַי אָנֹכִי מְבַקֵּשׁ הַגִּידָה נָּא לִי אֵיפֹה הֵם רֹעִים׃}
}
{
	\eE{
		And he said, “I’m looking for my brothers. Tell me, where
		are they feeding?”
	}
}
{
	He's like ``I'm looking for my brothers. I figured they might
	be out here in this here field. Have you seen them.''
}

\Verse{17}
{
	\hE{וַיֹּאמֶר הָאִישׁ נָסְעוּ מִזֶּה כִּי שָׁמַעְתִּי אֹמְרִים נֵלְכָה דֹּתָיְנָה וַיֵּלֶךְ יוֹסֵף אַחַר אֶחָיו וַיִּמְצָאֵם בְּדֹתָן׃}
}
{
	\eE{
		And the man said, “They traveled on from here. Because I
		heard them saying, ‘Let’s go to Dothan.’” And Joseph went
		after his brothers and found them in Dothan.
	}
}
{
	So the man's like ``Oh yeah those kids. I know the ones.
	They said they were going to Twin Pits.''
	
	So Joseph goes to Twin Pits and finds them.
}

\Verse{18}
{
	\hE{וַיִּרְאוּ אֹתוֹ מֵרָחֹק וּבְטֶרֶם יִקְרַב אֲלֵיהֶם וַיִּתְנַכְּלוּ אֹתוֹ לַהֲמִיתוֹ׃}
}
{
	\eE{
		And they saw him from a distance, and before he came close
		to them they conspired against him: to kill him.
	}
}
{
	They see him coming and they're like ``Ugh, Joe's here. Let's kill him.''
}

\Verse{19}
{
	\hJ{וַיֹּאמְרוּ אִישׁ אֶל אָחִיו הִנֵּה בַּעַל הַחֲלֹמוֹת הַלָּזֶה בָּא׃}
}
{
	\eJ{
		And the brothers said to one another, “Here comes the
		dream-master, that one there!
	}
}
{
	One of those brothers is like ``Hey the dream boy's coming.''
}

\Verse{20}
{
	\hJ{וְעַתָּה לְכוּ וְנַהַרְגֵהוּ וְנַשְׁלִכֵהוּ בְּאַחַד הַבֹּרוֹת וְאָמַרְנוּ חַיָּה רָעָה אֲכָלָתְהוּ וְנִרְאֶה מַה יִּהְיוּ חֲלֹמֹתָיו׃}
}
{
	\eJ{
		And now, come on and let’s kill him and throw him in one
		of the pits, and we’ll say a wild animal ate him, and
		we’ll see what his dreams will be!”
	}
}
{
	``Let's kill him and throw him in one of the pits. They have two
	in this town after all. We'll say a wild animal ate him. Bet he'll
	have really great dreams when he's dead ha!.''
%	\footnote{
%		wild animal. Literally, a bad animal—not just any wild animal. This
%		expression will occur one more time in the Torah: the blessings of the
%		covenant at Sinai will include a divine promise to eliminate wild (bad)
%		animals from the land. The value of that blessing is rendered more
%		visible by the recollection of Jacob’s pain when he believes that such a
%		creature has taken the life of his child.
%	}
}

\Verse{21}
{
	\hE{וַיִּשְׁמַע רְאוּבֵן וַיַּצִּלֵהוּ מִיָּדָם וַיֹּאמֶר לֹא נַכֶּנּוּ נָפֶשׁ׃}
}
{
	\eE{
		And Reuben heard, and he saved him from their hand. And he
		said, “Let’s not take his life.”
	}
}
{
	The eldest brother hears and he's like ``Guys c'mon, we're not
	gonna kill Joseph.''
}

\Verse{22}
{
	\hE{וַיֹּאמֶר אֲלֵהֶם רְאוּבֵן אַל תִּשְׁפְּכוּ דָם הַשְׁלִיכוּ אֹתוֹ אֶל הַבּוֹר הַזֶּה אֲשֶׁר בַּמִּדְבָּר וְיָד אַל תִּשְׁלְחוּ בוֹ לְמַעַן הַצִּיל אֹתוֹ מִיָּדָם לַהֲשִׁיבוֹ אֶל אָבִיו׃}
}
{
	\eE{
		And Reuben said to them, “Don’t spill blood. Throw him
		into this pit that’s in the wilderness, and don’t put out
		a hand against him”—in order to save him from their hand,
		to bring him back to his father.
	}
}
{
	``Just throw him in this pit.''
}

\Verse{23}
{
	\hJ{וַיְהִי כַּאֲשֶׁר בָּא יוֹסֵף אֶל אֶחָיו וַיַּפְשִׁיטוּ אֶת יוֹסֵף אֶת כֻּתּנְתּוֹ אֶת כְּתֹנֶת הַפַּסִּים אֲשֶׁר עָלָיו׃}
}
{
	\eJ{
		And it was when Joseph came to his brothers: and they took
		off Joseph’s coat, the coat of many colors, which he had
		on.
	}
}
{
	\aB{So Joseph catches up to the brothers and they
		steal his rad shirt.}
}

\Verse{24}
{
	\hE{וַיִּקָּחֻהוּ וַיַּשְׁלִכוּ אֹתוֹ הַבֹּרָה וְהַבּוֹר רֵק אֵין בּוֹ מָיִם׃}
}
{
	\eE{
		And they took him and threw him into the pit. And the pit
		was empty; there was no water in it.
	}
}
{
	And they throw him in the pit.
	
	It's ok, there's no water in there.

	He doesn't need water.
%	\footnote{
%		the pit was empty; there was no water in it. Rashi takes this to be
%		redundant: if it is empty, of course there is no water in it. He
%		concludes that there are snakes and scorpions in it. But it is not
%		redundant. Rather, two things are conveyed: It is empty. This conveys
%		that he is alone and helpless. There is no water in it. This conveys
%		that his survival is in danger.
%	}
}

\Verse{25}
{
	\hE{וַיֵּשְׁבוּ לֶאֱכל לֶחֶם}
	\hJ{וַיִּשְׂאוּ עֵינֵיהֶם וַיִּרְאוּ וְהִנֵּה אֹרְחַת יִשְׁמְעֵאלִים בָּאָה מִגִּלְעָד וּגְמַלֵּיהֶם נֹשְׂאִים נְכֹאת וּצְרִי וָלֹט הוֹלְכִים לְהוֹרִיד מִצְרָיְמָה׃}
}
{
	\eE{And they sat down to eat bread.}
	\eJ{And they raised their eyes and saw, and here was a caravan
		of Ishmaelites coming from Gilead, and their camels were carrying
		spices and balsam and myrrh, going to bring them down to Egypt.}
}
{
	Ok so remember how Hagar had a son before YHWH got Sarah pregnant
	and how his name was Ishmael?
	
	Apparently that guy's descendants have become a whole race at
	this point.\aB{See man, YHWH finally kept one of those seed promises.}

	So they sit down to eat and they see some Ishmaelites
	coming from Heap Of Testimony (hereafter HOT).
	
	Also their camels have like seasoning and stuff.
	
	They're headed to Egypt.
	
	Ishmael's mom was from Egypt remember?
}

\Verse{26}
{
	\hJ{וַיֹּאמֶר יְהוּדָה אֶל אֶחָיו מַה בֶּצַע כִּי נַהֲרֹג אֶת אָחִינוּ וְכִסִּינוּ אֶת דָּמוֹ׃}
}
{
	\eJ{
		And Judah said to his brothers, “What profit is there if
		we kill our brother and cover his blood?
	}
}
{
	So then Judah, who the country is named after later in the book,
	says one of the best lines in the bible.
	
	``Wait guys, we don't make any money if we kill him.''
%	\footnote{
%		What profit is there. Judah’s motives are a mystery. Is he really
%		concerned with money (“let’s sell him to the Ishmaelites”) and avoiding
%		performing fratricide (“let our hand not be on him, because he’s our
%		brother”)? Or is his motive to save Joseph’s life? Judah’s later
%		behavior toward Benjamin and his great reward at the end (see the
%		comment on 49:8) suggest that his behavior here is positive. Reuben’s
%		failed plan leads Joseph to the pit; Judah’s plan leads him to Egypt.
%		And Abraham’s unions with Hagar and Keturah produced the Ishmaelites and
%		the Midianites, who are now crucial to Joseph’s going from the pit to
%		Egypt! Perhaps the point is how complex and fragile our fate is.
%	}
}

\Verse{27}
{
	\hJ{לְכוּ וְנִמְכְּרֶנּוּ לַיִּשְׁמְעֵאלִים וְיָדֵנוּ אַל תְּהִי בוֹ כִּי אָחִינוּ בְשָׂרֵנוּ הוּא וַיִּשְׁמְעוּ אֶחָיו׃}
}
{
	\eJ{
		Come on and let’s sell him to the Ishmaelites, and let our
		hand not be on him, because he’s our brother, our flesh.”
		And his brothers listened.
	}
}
{
	Judah's like ``Look guys, I know we can't stand him, but
	he's our own flesh and blood, he's our brother,
	let's do the right thing and sell him into slavery.''

	The brothers are like ``Ok!''
}

\Verse{28}
{
	\hE{וַיַּעַבְרוּ אֲנָשִׁים מִדְיָנִים סֹחֲרִים וַיִּמְשְׁכוּ וַיַּעֲלוּ אֶת יוֹסֵף מִן הַבּוֹר}
	\hJ{וַיִּמְכְּרוּ אֶת יוֹסֵף לַיִּשְׁמְעֵאלִים בְּעֶשְׂרִים כָּסֶף וַיָּבִיאוּ אֶת יוֹסֵף מִצְרָיְמָה׃}
}
{
	\eE{And Midianite people, merchants, passed, and they pulled
		and lifted Joseph from the pit.}
	\eJ{And they sold Joseph to the Ishmaelites for twenty weights
		of silver. And they brought Joseph to Egypt.}
}
{
	\aB{Then some Midianites take him out of the pit.}

	And they\fB{Who?} sold Joseph to the Ishmaelites for twenty units
	of money.

	They bring him to Egypt.
}

\Verse{29}
{
	\hE{וַיָּשׁב רְאוּבֵן אֶל הַבּוֹר וְהִנֵּה אֵין יוֹסֵף בַּבּוֹר וַיִּקְרַע אֶת בְּגָדָיו׃}
}
{
	\eE{
		And Reuben came back to the pit, and here: Joseph was not
		in the pit. And he tore his clothes.
	}
}

{
       Reuben comes back to the pit and Josephs gone.

       He (presumably Reuben) tears his (presumably Reuben's) clothes.%
       \fB{It could also be the rad shirt.}
}

\Verse{30}
{
	\hE{וַיָּשׁב אֶל אֶחָיו וַיֹּאמַר הַיֶּלֶד אֵינֶנּוּ וַאֲנִי אָנָה אֲנִי בָא׃}
}
{
	\eE{
		And he went back to his brothers and said, “The boy’s
		gone! And I, where can I go?”
	}
}
{
       He goes back to the others and he's like ``Joseph's gone!
       What are we gonna tell dad?''
}

\Verse{31}
{
	\hJ{וַיִּקְחוּ אֶת כְּתֹנֶת יוֹסֵף וַיִּשְׁחֲטוּ שְׂעִיר עִזִּים וַיִּטְבְּלוּ אֶת הַכֻּתֹּנֶת בַּדָּם׃}
}
{
	\eJ{
		And they took Joseph’s coat and slaughtered a he-goat and
		dipped the coat in the blood.
	}
}
{
       So they take Joseph's rad shirt and dip it in goat blood.
%	\footnote{
%		they took Joseph’s coat and slaughtered a he-goat. Jacob was ironically
%		paid back for appropriating his brother’s birthright: he had to work an
%		additional seven years for Rachel on account of her sister’s birthright
%		(see the comment on Gen 29:26). Now he is likewise paid back for
%		appropriating his brother’s blessing: He deceived his father using his
%		brother’s clothing and the meat and skins of a goat. Now his own sons
%		deceive him using their brother’s clothing and the blood of a goat. The
%		Torah does not excuse Jacob’s behavior. It rather teaches that such acts
%		have consequences.
%	}
}

\Verse{32}
{
	\hJ{וַיְשַׁלְּחוּ אֶת כְּתֹנֶת הַפַּסִּים וַיָּבִיאוּ אֶל אֲבִיהֶם וַיֹּאמְרוּ זֹאת מָצָאנוּ הַכֶּר נָא הַכְּתֹנֶת בִּנְךָ הִוא אִם לֹא׃}
}
{
	\eJ{
		And they sent the coat of many colors and brought it to
		their father and said, “We found this. Recognize: is it
		your son’s coat or not?”
	}
}
{
	Then they deceive their dad (Jacob) using:
	
	\begin{enumerate}
	\item Their brother's clothes.
	\item A dead goat.
	\end{enumerate}

	Which coincidentally Jacob did to \textsl{his} dad!

	\aB{So they're like ``Dad, this shirt looks pretty rad. Is this Joseph's?''}
}

\Verse{33}
{
	\hJ{וַיַּכִּירָהּ וַיֹּאמֶר כְּתֹנֶת בְּנִי חַיָּה רָעָה אֲכָלָתְהוּ טָרֹף טֹרַף יוֹסֵף׃}
}
{
	\eJ{
		And he recognized it and said, “My son’s coat. A wild
		animal ate him. Joseph is torn up!”
	}
}
{
	So without the brothers \textsl{technically} lying, Jacob
	immediately comes to the conclusion they hoped he would.
	
	Joseph's been all eaten up by some animal.
}

\Verse{34}
{
	\hJ{וַיִּקְרַע יַעֲקֹב שִׂמְלֹתָיו וַיָּשֶׂם שַׂק בְּמתְנָיו וַיִּתְאַבֵּל עַל בְּנוֹ יָמִים רַבִּים׃}
}
{
	\eJ{
		And Jacob ripped his clothes and wore sackcloth on his
		hips and mourned over his son many days.
	}
}
{
	\aB{Dad does some behaviors that seem a bit wacky to us these days
		but they're totally standard mourning practices in the Ancient
		Near East.}
}

\Verse{35}
{
	\hJ{וַיָּקֻמוּ כל בָּנָיו וְכל בְּנֹתָיו לְנַחֲמוֹ וַיְמָאֵן לְהִתְנַחֵם וַיֹּאמֶר כִּי אֵרֵד אֶל בְּנִי אָבֵל שְׁאֹלָה וַיֵּבְךְּ אֹתוֹ אָבִיו׃}
}
{
	\eJ{
		And all his sons and all his daughters got up to console
		him, and he refused to be consoled, and he said, “Because
		I’ll go down mourning to my son at Sheol,” and his father
		wept for him.
	}
}
{
	All his sons%
	\footnote{Eleven.} and daughters%
	\footnote{
		What's left of Dinah, but plural for some reason.
	} try to console him.

	But he's sad anyway.
	
	Cuz obviously.
	
	Joseph was Rachel's
	
	Now he has nothing left of Rachel except stupid Benjamin who killed her
	in childbirth.
}

\Verse{36}
{
	\hE{וְהַמְּדָנִים מָכְרוּ אֹתוֹ אֶל מִצְרָיִם לְפוֹטִיפַר סְרִיס פַּרְעֹה שַׂר הַטַּבָּחִים׃}
}
{
	\eE{
		And the Medanites sold him to Egypt, to Potiphar, an
		official of Pharaoh, chief of the guards.
	}
}
{
	So now the MDINIM (Heb. \heb{מִדְיָנִים}) sell him to Egypt.
	
	They're probably the same as the MDNIM (\heb{מְּדָנִים}),
	but spelled different for some reason.

	Anyways they sell him to an Egyptian named Potiphar,
	who we'll hear more about in the next chapter.
}
